\chapter{一致可积性}
\label{chap:uniform-integrability}

控制收敛定理用一个可积函数同时支配整列。许多函数族没有共同点态上界，
但它们在大函数值区域所携带的积分质量仍可统一变小。一致可积性把这种性质
直接写成积分条件。它与依测度收敛结合，给出 \(L^1\) 收敛；在概率论和弱
紧致理论中，它常比寻找控制函数更合适。

\section{小集合与高函数值}

本章固定测度空间 \((X,\mathcal A,\mu)\)，函数均为实可测函数。
可积函数取有限可测代表，积分条件不依赖代表的选择。数列指标从
\(n=1\) 开始。设 \(\mathcal F\subset L^1(\mu)\)；空函数族的非负积分
上确界约定为零。考虑两种控制，其中 \(E\in\mathcal A\)：
\[
\tag{UI-tail}
\lim_{M\to\infty}\sup_{f\in\mathcal F}
\int_{\{|f|>M\}}|f|\,d\mu=0,
\]
以及
\[
\tag{UI-set}
\forall\varepsilon>0\ \exists\delta>0:\quad
\mu(E)<\delta\Longrightarrow
\sup_{f\in\mathcal F}\int_E|f|\,d\mu<\varepsilon.
\]
第一式排除越来越高的尖峰；第二式说没有函数能在很小的集合上集中固定质量。
对单个 \(L^1\) 函数，二者都由积分绝对连续性得到；对函数族，量词中的
\(\sup_{f\in\mathcal F}\) 是新增要求。

\begin{definition}
若 \(\mathcal F\subset L^1(\mu)\) 满足 \textup{(UI-tail)}，则称
\(\mathcal F\) \term{一致可积}。
\end{definition}

\begin{proposition}[有限测度空间上的判据]
\label{prop:ui-finite-measure-criterion}
若 \(\mu(X)<\infty\)，则 \(\mathcal F\) 一致可积当且仅当：
\begin{enumerate}[label=\textup{(\roman*)}]
 \item \(\sup_{f\in\mathcal F}\|f\|_1<\infty\)；
 \item \(\mathcal F\) 满足 \textup{(UI-set)}。
\end{enumerate}
\end{proposition}

\begin{proof}
先设一致可积。取 \(M>0\) 使尾积分统一小于 \(1\)，则
\[
 \int|f|\le M\mu(X)+1,
\]
得到 (i)。给定 \(\varepsilon>0\)，取 \(M>0\) 使尾积分小于
\(\varepsilon/2\)，再取 \(\delta=\varepsilon/(2M)\)。于是
\[
 \int_E|f|
 \le M\mu(E)+\int_{\{|f|>M\}}|f|<\varepsilon.
\]

反过来，设 \(C=\sup_f\|f\|_1<\infty\)。由 Markov 不等式，
\[
 \mu(\{|f|>M\})\le C/M.
\]
给定 \(\varepsilon>0\)，先由 (ii) 取 \(\delta\)，再令
\(M>\max\{C/\delta,0\}\)。将 \(E=\{|f|>M\}\) 代入 (ii)，便得尾积分统一小于
\(\varepsilon\)。
\end{proof}

\section{例子与反例}

上述证明中，有限测度只用于从尾积分控制推出整体 \(L^1\) 有界性。
因此，在任意测度空间上，\(L^1\) 有界加上 \textup{(UI-set)} 仍足以推出
一致可积；一致可积也总能推出 \textup{(UI-set)}。不能省略的是反向所需的
整体界：在 \(\mathbb R\) 上，\(\{\mathbf1_{[0,n]}:n\ge1\}\) 的尾积分在
\(M>1\) 时全为零，但各函数的积分等于 \(n\)。本章采用的尾积分定义因而
不在无限测度空间上自动包含 \(L^1\) 有界性。

\begin{proposition}
\label{prop:dominated-family-ui}
若存在 \(g\in L^1(\mu)\) 使 \(|f|\le g\) 对所有
\(f\in\mathcal F\) 几乎处处成立，则 \(\mathcal F\) 一致可积。
\end{proposition}

\begin{proof}
在 \(\{|f|>M\}\) 上必有 \(g>M\)，故
\[
 \int_{\{|f|>M\}}|f|
 \le\int_{\{g>M\}}g\longrightarrow0.
\]
\end{proof}

\begin{counterexample}[尖峰]
在 \((0,1)\) 上，
\[
 f_n=n\mathbf1_{(0,1/n)}
\]
满足 \(\|f_n\|_1=1\)，但不一致可积，因为对任意 \(M\)，只要
\(n>M\)，便有
\[
 \int_{\{|f_n|>M\}}|f_n|=1.
\]
\end{counterexample}

\begin{counterexample}[质量逃逸]
在 \(\mathbb R\) 上，
\[
 f_n=\mathbf1_{[n,n+1]}
\]
是一致可积的：当 \(M>1\) 时所有尾积分为零。它还满足小集合控制和
\(L^1\) 有界，却在局部依测度趋于零而不在 \(L^1\) 中趋于零。问题不是
高函数值，而是质量移向任何固定有限区域之外。
\end{counterexample}

\section{凸函数控制}

\begin{theorem}[de la Vallée--Poussin 判据]
\label{thm:de-la-vallee-poussin}
若存在递增凸函数 \(\Phi:[0,\infty)\to[0,\infty)\)，满足
\[
 \frac{\Phi(t)}t\longrightarrow\infty,
 \qquad
 \sup_{f\in\mathcal F}\int_X\Phi(|f|)\,d\mu<\infty,
\]
则 \(\mathcal F\) 一致可积。反过来，在有限测度空间上，一致可积族可以
找到这样的分段线性凸函数 \(\Phi\)。
\end{theorem}

\begin{proof}
记 \(C=\sup_f\int\Phi(|f|)\)。给定 \(\varepsilon>0\)，由
\(\Phi(t)/t\to\infty\)，可取 \(M\) 使
\[
 t\le\frac{\varepsilon}{C+1}\Phi(t)\qquad(t>M).
\]
于是
\[
 \sup_f\int_{\{|f|>M\}}|f|
 \le\frac{\varepsilon}{C+1}\sup_f\int\Phi(|f|)
 \le\varepsilon\frac C{C+1}<\varepsilon.
\]

反过来，设 \(\mu(X)<\infty\) 且 \(\mathcal F\) 一致可积。由命题
\ref{prop:ui-finite-measure-criterion}，存在
\[
 C_0=\sup_{f\in\mathcal F}\|f\|_1<\infty.
\]
逐次选取严格递增的 \(M_k\to\infty\)，使
\[
 \sup_f\int_{\{|f|>M_k\}}|f|<2^{-k}.
\]
定义
\[
 \Phi(t)=t+\sum_{k=1}^{\infty}(t-M_k)_+,
 \qquad t\ge0.
\]
对每个固定 \(t\)，和式中只有有限项非零。函数 \(\Phi\) 连续、递增、分段
线性，斜率在经过每个 \(M_k\) 时增加 \(1\)，因而凸。对任意固定的
\(N\)，当 \(t\ge2\max_{k\le N}M_k\) 时，
\[
 \frac{\Phi(t)}t
 \ge\sum_{k=1}^{N}\frac{t-M_k}{t}\ge\frac N2.
\]
所以 \(\Phi(t)/t\to\infty\)。最后，Tonelli 定理给出
\[
 \begin{split}
 \int\Phi(|f|)
 &=\int|f|+\sum_{k=1}^{\infty}\int(|f|-M_k)_+\\
 &\le C_0+\sum_{k=1}^{\infty}
       \int_{\{|f|>M_k\}}|f|
 \le C_0+1.
 \end{split}
\]
右端与 \(f\in\mathcal F\) 无关，反向构造完成。
\end{proof}

\begin{example}
若对某个 \(p>1\)，
\[
 \sup_{f\in\mathcal F}\|f\|_p<\infty,
\]
则取 \(\Phi(t)=t^p\) 得 \(\mathcal F\) 一致可积。因而有限测度空间上的
\(L^p\) 有界族（\(p>1\)）在 \(L^1\) 意义下不会形成高尖峰。\(p=1\)
时结论正是尖峰反例所否定的内容。
\end{example}

\begin{example}[没有共同可积控制函数的一致可积列]
在 \((0,1)\) 上取互不相交的区间
\[
 I_n=\left(\frac1{2(n+1)},\frac1{2n}\right),\qquad
 f_n=n\mathbf1_{I_n}.
\]
由于
\[
 \int_0^1 |f_n|^2
 =n^2\left(\frac1{2n}-\frac1{2(n+1)}\right)
 =\frac{n}{2(n+1)}\le\frac12,
\]
前例说明这列函数一致可积。若有 \(g\in L^1(0,1)\) 对每个 \(n\)
几乎处处支配 \(|f_n|\)，删去相应零集的可数并后，在每个 \(I_n\)
上均有 \(g\ge n\)。利用区间互不相交和积分的可数可加性，
\[
 \int_0^1 |g|\ge\sum_{n=1}^{\infty}n\,m(I_n)
 =\sum_{n=1}^{\infty}\frac1{2(n+1)}=\infty,
\]
矛盾。这说明寻找一个共同控制函数可能失败，而一致可积判据仍可适用。
\end{example}

\section{Vitali 收敛定理}

\begin{theorem}[Vitali]
\label{thm:vitali-convergence}
设 \(\mu(X)<\infty\)，\(f_n\in L^1(\mu)\)，且
\(f_n\xrightarrow{\mu}f\)。若 \(\{f_n:n\ge1\}\) 一致可积，则
\(f\in L^1\)，并且
\[
 \|f_n-f\|_1\longrightarrow0.
\]
反过来，\(L^1\) 收敛蕴含依测度收敛，并且
\(\{f_n:n\ge1\}\cup\{f\}\) 一致可积。
\end{theorem}

\begin{proof}
先由子列原理抽出 \(f_{n_k}\to f\) a.e.。Fatou 引理和一致的
\(L^1\) 有界性给出
\(\int|f|\le\liminf_k\int|f_{n_k}|<\infty\)。

给定 \(\varepsilon>0\)。一致可积性及 \(f\in L^1\) 的积分绝对连续性
给出 \(\delta>0\)，使 \(\mu(E)<\delta\) 时
\[
 \sup_n\int_E|f_n|<\varepsilon,\qquad
 \int_E|f|<\varepsilon.
\]
选 \(\eta>0\) 使 \(\eta\mu(X)<\varepsilon\)。依测度收敛给出
\[
 E_n=\{|f_n-f|>\eta\},\qquad \mu(E_n)<\delta
\]
最终成立。于是
\[
 \int|f_n-f|
 \le\eta\mu(X)+\int_{E_n}|f_n|+\int_{E_n}|f|
 <3\varepsilon.
\]

反向结论中，依测度收敛由 Markov 不等式。为证一致可积，固定
\(\varepsilon>0\)，先取 \(N\) 使 \(n\ge N\) 时
\(\|f_n-f\|_1<\varepsilon\)。有限多个
\(f_1,\ldots,f_{N-1},f\) 的积分具有共同的绝对连续性；再用
\(\int_E|f_n|\le\int_E|f|+\|f_n-f\|_1\) 得小集合控制和统一
\(L^1\) 有界性，应用命题 \ref{prop:ui-finite-measure-criterion}。
\end{proof}

\begin{corollary}
在有限测度空间上，控制收敛定理是 Vitali 定理的特例：若
\(|f_n|\le g\in L^1\)，则 \(\{f_n\}\) 一致可积，而几乎处处收敛蕴含
依测度收敛。一般测度空间上的控制收敛定理仍由第
\ref{chap:lebesgue-convergence} 章的直接证明给出，不需要有限测度假设。
\end{corollary}

\section{无限测度空间上的质量紧性}

在无限测度空间中，一致可积性只排除高尖峰，不能阻止积分质量移向越来越远
的区域。还需要一个控制空间尾部的条件。

\begin{definition}
设 \((X,\mathcal A,\mu)\) 为 \(\sigma\)-有限空间。若
\(\mathcal F\subset L^1(\mu)\) 满足：对每个 \(\varepsilon>0\)，存在
\(K\in\mathcal A\)、\(\mu(K)<\infty\)，使
\[
 \sup_{f\in\mathcal F}\int_{X\setminus K}|f|\,d\mu<\varepsilon,
\]
则称 \(\mathcal F\) 具有\term{质量紧性}。
\end{definition}

这里的 \(K\) 只要求可测且测度有限；若空间带有拓扑并具有正则测度，在常见
情形下还可进一步选取紧集。术语中的“紧”描述积分质量不能逃出所有有限区域，
不应与集合本身的拓扑紧致性混同。

\begin{theorem}[无限测度空间上的 Vitali 定理]
\label{thm:vitali-sigma-finite-tight}
设 \((X,\mathcal A,\mu)\) 为 \(\sigma\)-有限测度空间，
\(f_n\in L^1(\mu)\)，且 \(f\) 是几乎处处有限的可测函数。假定：
\begin{enumerate}[label=\textup{(\roman*)}]
 \item 对每个有限测度可测集 \(E\)，\(f_n\to f\) 在 \(E\) 上依测度收敛；
 \item 函数族 \(\{f_n:n\ge1\}\) 一致可积；
 \item 该函数族具有质量紧性。
\end{enumerate}
则 \(f\in L^1(\mu)\)，并且
\[
 \|f_n-f\|_1\longrightarrow0.
\]
\end{theorem}

\begin{proof}
先构造一个在全空间几乎处处收敛的子列。取有限测度集
\(X_j\uparrow X\)。由 (i)，可递归选取严格递增的指标 \(n_j\)，使
\[
 \mu\bigl(X_j\cap\{|f_{n_j}-f|>2^{-j}\}\bigr)<2^{-j}.
\]
这些坏集的测度和有限，故其上极限为零测集。对不属于该上极限且属于
\(\bigcup_jX_j=X\) 的点，最终有
\(|f_{n_j}-f|\le2^{-j}\)；所以 \(f_{n_j}\to f\) a.e.

由 (iii)，可取有限测度集 \(K\)，使
\(\sup_n\int_{X\setminus K}|f_n|<1\)。把 (ii) 限制到 \(K\) 上仍得到
一致可积族，而 (i) 给出 \(K\) 上的依测度收敛。有限测度空间上的 Vitali
定理说明 \(f\in L^1(K)\)。另一方面，对刚才的几乎处处收敛子列使用 Fatou
引理，得到
\[
 \int_{X\setminus K}|f|
 \le\liminf_{j\to\infty}
     \int_{X\setminus K}|f_{n_j}|
 \le1.
\]
故 \(f\in L^1(X)\)。

最后给定 \(\varepsilon>0\)，再由质量紧性取有限测度集 \(K\)，使
\[
 \sup_n\int_{X\setminus K}|f_n|<\varepsilon.
\]
同一 Fatou 论证给出
\(\int_{X\setminus K}|f|\le\varepsilon\)，而有限测度 Vitali 定理给出
\(\int_K|f_n-f|\to0\)。因此
\[
 \limsup_{n\to\infty}\|f_n-f\|_1
 \le 2\varepsilon.
\]
令 \(\varepsilon\downarrow0\)，得到结论。
\end{proof}

移动指示函数 \(\mathbf1_{[n,n+1]}\) 局部依测度趋于零且一致可积，但它的
单位积分质量离开每个固定有限测度区域，因而不满足质量紧性。

\section*{习题}
\addcontentsline{toc}{section}{习题}
\subsection*{A组　判据与例子}
\begin{enumerate}[label=\textbf{\thechapter.\arabic*.}]
\exerciseitem{ui-equivalent-truncation}
证明一致可积等价于
\[
 \lim_{M\to\infty}\sup_{f\in\mathcal F}
 \int\bigl(|f|-M\bigr)_+\,d\mu=0.
\]

\exerciseitem{ui-finite-family}
证明 \(L^1\) 中任意有限函数族一致可积。

\exerciseitem{ui-dominated-family}
用积分绝对连续性而不是尾积分，重新证明由同一 \(g\in L^1\) 控制的函数族
一致可积。

\exerciseitem{ui-spike-details}
验证 \(n\mathbf1_{(0,1/n)}\) 满足 \(L^1\) 有界和依测度收敛，却不一致
可积。

\end{enumerate}
\subsection*{B组　Vitali 定理}
\begin{enumerate}[label=\textbf{\thechapter.\arabic*.},resume]
\exerciseitem{ui-lp-bounded}
若 \(p>1\)、\(M>0\) 且 \(\sup_f\|f\|_p\le C<\infty\)，证明
\[
 \int_{\{|f|>M\}}|f|\le C^pM^{1-p}.
\]

\exerciseitem{ui-convex-sufficient}
补全 de la Vallée--Poussin 判据充分性证明，并说明凸性在这一方向是否
实际用到。

\exerciseitem{vitali-ae-version}
在有限测度空间上，证明若 \(f_n\to f\) a.e. 且 \(\{f_n\}\) 一致可积，
则 \(f_n\to f\) 于 \(L^1\)。

\exerciseitem{l1-convergence-ui}
不用引用 Vitali 定理的反向部分，直接证明 \(L^1\) 收敛列连同其极限构成
一致可积族。

\end{enumerate}
\subsection*{C组　紧性与推广}
\begin{enumerate}[label=\textbf{\thechapter.\arabic*.},resume]
\exerciseitem{ui-sums}
若 \(\mathcal F,\mathcal G\) 一致可积，证明
\(\mathcal F+\mathcal G=\{f+g:f\in\mathcal F,g\in\mathcal G\}\)
一致可积。

\exerciseitem{mass-tightness-translation}
证明 \(\{\mathbf1_{[n,n+1]}:n\ge1\}\) 一致可积但不满足质量紧性。

\optionalexerciseitem{vitali-infinite-space}
重作定理 \ref{thm:vitali-sigma-finite-tight} 的证明，尤其写清对角子列怎样
在全空间几乎处处收敛，并分别说明三个假设排除了哪一类反例。

\projectexerciseitem{ui-random-variables}
总测度为1的测度空间称为概率空间，实可测函数称为随机变量，其积分记为
期望 \(\mathbb E\)，依测度收敛在这里称为依概率收敛。
设 \(0<p<\infty\)。把一致可积性写成期望条件，并证明：若实随机变量
\(X_n\to X\) 依概率且 \(\{|X_n|^p:n\ge1\}\) 一致可积，则
\(X_n,X\in L^p\) 且 \(\mathbb E|X_n-X|^p\to0\)。
\end{enumerate}
